\documentclass[10pt,a4paper]{amsart}
\usepackage[margin=19mm]{geometry}
\usepackage{amsmath,amssymb,mathtools}
\usepackage[scr=rsfs,cal=euler]{mathalfa}
\usepackage{xfrac}
\usepackage[unicode,hidelinks,bookmarks=false]{hyperref}
\newcommand{\ie}{i.e.\ }
\newcommand{\cf}{cf.\ }
\newcommand{\resp}{respectively\ }
\newcommand{\Subsection}{\subsection}
\providecommand{\nobreakdash}{\nobreak\hbox{-}}
\setlength{\emergencystretch}{2em}
\multlinegap0pt
\usepackage{bm}
\usepackage{bbm}
\usepackage{dsfont}
\usepackage{xspace}
\usepackage{cancel}
\usepackage{mathtools}
\usepackage{amsthm}
\makeatletter
\@ifundefined{theorem}{\newtheorem{theorem}{Theorem}}{}
\makeatother
\title[On a novel approach to nonexpansive mappings]{On a novel approach to nonexpansive mappings}
\author{Anish Banerjee, Hiranmoy Garai, Pratikshan Mondal, Lakshmi Kanta Dey}
\date{Source: arXiv:2508.06821v1 (CC BY 4.0)}
\begin{document}
\maketitle

\noindent\emph{Source and licence.} On a novel approach to nonexpansive mappings. Version arXiv:2508.06821v1 (CC BY 4.0), distributed under the Creative Commons Attribution 4.0 International licence, as recorded in the arXiv licence metadata for this exact version and shown on the abstract page https://arxiv.org/abs/2508.06821v1. Full rights notice: licenses/arxiv-2508.06821-CC-BY-4.0.txt.\par\smallskip

\noindent\textit{Editorial scope.} Counted unit: the Theorem whose source label is \texttt{T35} in the article (source0.tex lines 312--314, source label \texttt{T35}), delivered together with the article's own complete original proof of it (source0.tex lines 316--352, one \texttt{proof} environment of 37 lines). No mathematics is written, repaired or re-derived: statement and proof are the article's own text line for line. Required context retained uncounted: none. Every definition, display and label of those lines is retained unchanged. Omitted from this package and not redistributed: 1--311, 315--315, 353--420. Nothing inside a retained excerpt is omitted or altered: every retained body is delivered exactly as the article's lines are written, so no omission receipt of any kind accompanies this package. Citations: the retained excerpts carry no citation command at all, so no bibliography entry of the article is transcribed and no reference is redistributed. Reference adaptation: none was needed and none was made; every reference inside the delivered unit resolves inside it, so no unresolved reference or citation is produced. Printed slips kept literal: no printed word, symbol, hypothesis or number of the counted statement or of its proof was altered. Layout and class: the article's own class is \texttt{amsart} with packages including amsmath, amsthm, amscd, amsfonts, amssymb, color, cite, graphicx, xcolor, eucal, upgreek, hyperref; the wrapper is the lane's pinned \texttt{amsart} editorial wrapper at 10pt on A4, which restates the theorem-like environments the retained text needs and adds the article's own macro definitions that the retained text needs (none), with extra wrapper packages (bm, bbm, dsfont, xspace, cancel, mathtools). The delivered numbering is the unit's own. Assets: no retained span contains a figure environment, a graphics-inclusion command, a PDF-page inclusion, a table or a TikZ picture; no image file is copied into this package, the package declares no asset dependency and no image file is redistributed with it. Licence: version arXiv:2508.06821v1 of this article is distributed under the Creative Commons Attribution 4.0 International licence, as recorded in the arXiv licence metadata for this exact version and shown on the abstract page https://arxiv.org/abs/2508.06821v1. A full rights notice is delivered at licenses/arxiv-2508.06821-CC-BY-4.0.txt. Characters: every retained excerpt is pure ASCII, so the delivered engine, XeLaTeX with its default Latin Modern text font, reports no missing character.
\begin{theorem}\label{T35}
Let $K$ be a closed, bounded, convex subset of the Hilbert space $H$ and  $T: K \to K$ be a perimetric nonexpansive mapping. Then $T$ has a fixed point.
\end{theorem}

\begin{proof}
Let $\{s_n\}$ be a sequence of positive real numbers satisfying $s_n < 1$ for each $n$, with the property that $\lim_{n \to \infty} s_n = 1$.
Let $x_0 \in K$  be fixed. Then since $K$ is convex, for each $n\in \mathbb{N}$, the mapping  $U_n: K \to K$ defined by 
\begin{align*}
U_n(x)= (1-s_n)x_0 +s_nTx, \text{ for all } x \in K
\end{align*} is well defined. Moreover, $U_n$ is a mapping contracting perimeter of triangles on $K$, and given that $K$ is closed, it follows that $U_n$ has a fixed point, say $u_n$ for each $n$.

Since $K$ is closed, convex and bounded in the Hilbert space $K$, it is weakly compact. Then we  get a subsequence $\{u_{n_k}\}$ of $\{u_n\}$ such that $\{u_{n_k}\}$ converges weakly to an element $p \in K$. Next, we show that $p$ is a fixed point of $T$.

For any  $u \in K$, we have 
\begin{align*}
{\lVert u_{n_k} - u \rVert}^2 
&= {\lVert u_{n_k} -p + p - u \rVert}^2\\
&= {\lVert u_{n_k} - p \rVert}^2 + {\lVert p - u \rVert}^2 + 2\langle u_{n_k} -p, p-u \rangle.
\end{align*}
Here $\langle u_{n_k} -p, p-u \rangle \to \theta$ as $k \to \infty$ as $u_{n_k} - p$ converges to $\theta$ weakly in $K$.
Then taking $u=Tp$, we have
\begin{align*}
\lim_{k \to \infty} ({\lVert u_{n_k} - Tp \rVert}^2 - {\lVert u_{n_k} - p \rVert}^2) = {\lVert Tp - p \rVert}^2.
\end{align*}
Since $s_{n_k} \to 1$ and $U_{n_k}u_{n_k}=u_{n_k}$, we have
\begin{align*}
Tu_{n_k} - u_{n_k} &= (s_{n_k} Tu_{n_k}+ (1-s_{n_k})x_0) - u_{n_k} + (1-s_{n_k})(Tu_{n_k} - x_0)\\
&=(U_{n_k}u_{n_k} -u_{n_k}) + (1-s_{n_k})(Tu_{n_k} - x_0)\\
&=\theta + (1- s_{n_k})(Tu_{n_k} - x_0). 
\end{align*}
Thus, we have $\lim_{k \to \infty} \lVert Tu_{n_k} - u_{n_k} \rVert =0$.
Finally, we obtain
\begin{align*}
&\lVert u_{n_k} - Tp \rVert \le \lVert u_{n_k} - Tu_{n_k} \rVert+ \lVert Tu_{n_k} - Tp \rVert + \lVert Tp - Tu_{n_{k+1}} \rVert+ \lVert Tu_{n_{k+1}} - Tu_{n_k} \rVert\\
\implies &\lVert u_{n_k} - Tp \rVert \le \lVert u_{n_k} - Tu_{n_k} \rVert+ \lVert u_{n_k} - p \rVert + \lVert p - u_{n_{k+1}} \rVert+ \lVert u_{n_{k+1}} - u_{n_k} \rVert\\
\implies & \lim_{k \to \infty} \lVert u_{n_k} - Tp \rVert \le \lim_{k \to \infty} \lVert u_{n_k} - Tu_{n_k} \rVert =0\\
\implies & \lVert p - Tp \rVert =0\\
\implies &Tp=p.
\end{align*}
Hence, $p$ is a fixed point of $T$.
\end{proof}

\end{document}
